% \zenodoCodeID is the one DOI the paper cites for the code, in the data
% availability section and in the entry BMB26 of the bibliography: the version
% DOI of release v5.0.1. Releases v1.0.0, v2.0.0, v3.0.0, v3.1.1, v3.2.0,
% v4.0.0 and v5.0.0 are 10.5281/zenodo.22950276, 10.5281/zenodo.23038895,
% 10.5281/zenodo.23047883, 10.5281/zenodo.23078541, 10.5281/zenodo.23093099,
% 10.5281/zenodo.23197107 and 10.5281/zenodo.23221435; all are versions of the
% concept record 10.5281/zenodo.22950275.
\newcommand{\zenodoCodeID}{10.5281/zenodo.23225483}
\newcommand{\zenodoCodeDOI}{%
  \href{https://doi.org/\zenodoCodeID}{\texttt{\zenodoCodeID}}}

\section*{Declarations}

\subsection*{Funding}
The authors received no funding for this work.

\subsection*{Competing interests}
The authors have no competing interests to declare.

\subsection*{Ethics approval and consent}
This article reports no research involving human participants, human data or
animals, so ethics approval and consent to participate or to publish were not
required. The results are original and have not been published elsewhere,
and the work of others is credited by citation throughout.

\subsection*{Data availability}
The programs of the archive \cite{BMB26}, which repeat much of the finite
arithmetic of the proofs as an independent check, the Lean~4 certificate of
part of that arithmetic, the Macaulay2
computations of local $\Ext$ algebras with their output, the sources of the
paper, and the generators and TikZ sources of every figure are openly
available on Zenodo under the DOI \zenodoCodeDOI; the version that
accompanies this paper is release 5.0.1 of the archive. A document in the
archive lists, for each result of the paper, the computations that check
it. The programs perform $2218$ checks, all of which pass, and the
Lean kernel checks four hundred and nine theorems, each depending on no axiom
or on propositional extensionality alone. Deep Bhattacharjee computed all
the calculations, in C, Python, Julia, Macaulay2, Lean~4 and shell scripts,
with the assistance of Claude Code, an AI model developed by Anthropic. The
programs in the archive are those in Python, which use the C libraries FLINT
and Arb through python-flint, in Macaulay2 and in shell scripts, together
with the Lean~4 certificate, in which he formalised the finite arithmetic.
Every result they certify can be reproduced from the archive
with the programs and the Lean kernel, and the authors are responsible for
the mathematical statements and their proofs. No other data were generated
or analysed.
